% RESULTADO_RECUPERADO: funtor dimensional integral, eq:cZ-internos,
% eq:longitud-ruta; composición reunida en Artículo V REV02, sección 69.
% Incorporación autónoma a III: mismo transporte, mismas hojas y misma carta.
\subsection{Constitutive speed and kinematic realization of the clock}
\label{iii:sec:enlace-velocidad-constitutiva}

The constitutive response and the kinematic readout realize the same
speed section. Their identification preserves the order of construction:
the APP sheets, their TRIT orientation, and the angular transport expressed
by TPK precede the speed readouts and the choice of units.
In the two-sheet representation of
Section~\ref{iii:sec:hojas}, with $x>|y|$ and rank-one projectors
$P_\pm$, the transport is
\[
 T=e^{-xI-y\mathcal R}=q_+P_++q_-P_-,
 \qquad q_\pm=e^{-x\mp y},\qquad 0<q_\pm<1.
\]
Its temporal blocks of lengths $90$ and $120$ produce the response
$\mathcal V$ in \eqref{iii:eq:vacuum-op}. The coordinates $x,y$ and the
channels are outputs of the preceding angular development; this composition
applies its readouts to that already constructed transport.

The symmetric logarithmic readout of temporal transport is
\begin{equation}
 \mathcal E=
 \log\!\bigl[(1-q_+^{90})(1-q_-^{90})\bigr]
 -\log\!\bigl[(1-q_+^{120})(1-q_-^{120})\bigr].
 \label{iii:eq:velocidad-lector-simetrico}
\end{equation}
On the dimensional speed line, with positive section $u_C$, the
kinematic readout expresses $c_{\mathrm{int}}=\exp(-\mathcal E)u_C$.
The constitutive response expresses
$c_{\mathrm{vac}}=\widehat c u_C$ through \eqref{iii:eq:four}.

\begin{proposition}[Constitutive and kinematic identity]
\label{iii:prop:identidad-velocidad}
In the same dimensional chart,
\begin{equation}
 c_{\mathrm{int}}=c_{\mathrm{vac}}=\widehat c\,u_C,
 \qquad \widehat c=(r_+r_-)^{-1}=\exp(-\mathcal E).
 \label{iii:eq:identidad-velocidad}
\end{equation}
The identification preserves sheet interchange.
\end{proposition}
\begin{proof}
All factors in \eqref{iii:eq:velocidad-lector-simetrico} are
positive. The logarithm law and \eqref{iii:eq:rchannels} give
\[
 \mathcal E=\log(r_+r_-)=\log\det\mathcal V,
 \qquad e^{-\mathcal E}=(\det\mathcal V)^{-1}.
\]
The determinant corresponds to the realization with two rank-one sheets.
Both definitions therefore express the same section. Sheet interchange
permutes $r_+$ and $r_-$ and preserves their product. The determinant
$\det T=e^{-2x}$ corresponds to elementary transport, whereas
$\det\mathcal V=r_+r_-$ corresponds to its temporal response;
\eqref{iii:eq:identidad-velocidad} uses the latter.
\end{proof}

\subsubsection{Path length and elementary duration}

The positive dimensional frame $(u_S,u_C,u_T,u_Q)$ induces the length
basis $u_L=u_Cu_T$. For an enriched path $\gamma$ with
$n=|\gamma|$ transitions and constant channel, the additive readouts are
\begin{equation}
 T_{\mathrm{int}}(\gamma)=nu_T,
 \qquad L_{\mathrm{int}}(\gamma)=n\widehat c\,u_L.
 \label{iii:eq:velocidad-longitud-reloj}
\end{equation}
Concatenation adds the numbers of transitions and their realizations.
Enriched routes also preserve their orientation and memory.

\begin{corollary}[Normalization of the elementary length]
For $n>0$, $L_{\mathrm{int}}(\gamma)/T_{\mathrm{int}}(\gamma)
=c_{\mathrm{vac}}$. In the normalization $t_0=u_T$,
\begin{equation}
 \ell_0=c_{\mathrm{int}}t_0=\widehat c\,u_L.
 \label{iii:eq:longitud-elemental-constitutiva}
\end{equation}
In a temporal chart $t_0=\tau_0u_T$, with $\tau_0>0$, one obtains
$\ell_0=\tau_0\widehat c\,u_L$.
\end{corollary}
\begin{proof}
Dividing the two expressions in \eqref{iii:eq:velocidad-longitud-reloj}
cancels $n$ and uses $u_L/u_T=u_C$. Multiplying the speed section
by $t_0$ proves both elementary formulas.
\end{proof}

When the channel depends on the edge, the readout retains the additive form
$L_{\mathrm{int}}(\gamma)=\sum_{e\in\gamma}\widehat c(e)u_L$.
Its quotient by $nu_T$ is the mean of the speeds of those edges.
The constant-channel expression is the homogeneous specialization of
this readout.

\subsubsection{Change of chart and adapted normalization}

\begin{proposition}[Covariance of the identification]
Let $u_C'=d u_C$ and $u_T'=\eta u_T$, with $d,\eta>0$.
The coordinates of the same sections satisfy
\begin{equation}
 \widehat c'=d^{-1}\widehat c,
 \qquad \tau_0'=\eta^{-1}\tau_0,
 \qquad u_L'=d\eta u_L.
 \label{iii:eq:velocidad-cambio-carta}
\end{equation}
If another representation writes $c'=\widehat c' u_C'$ with $u_C'=d u_C$,
equality of sections is exactly equivalent to $d\widehat c'=\widehat c$.
\end{proposition}
\begin{proof}
Equality of each section in the two bases gives
$\widehat c'u_C'=\widehat c u_C$ and $\tau_0'u_T'=\tau_0u_T$.
In particular,
\[
 \tau_0'\widehat c'u_L'
 =\frac{\tau_0}{\eta}\frac{\widehat c}{d}\,d\eta u_L
 =\tau_0\widehat c u_L=\ell_0.
\]
The final equivalence follows from the same substitution for speed.
\end{proof}

The adapted choice $d=\widehat c$ gives $u_C'=c_{\mathrm{vac}}$ and
$\widehat c'=1$: it expresses the constructed section in an adapted basis.
The coefficient in the internal chart remains $(r_+r_-)^{-1}$, and that
is the normalization used by cubic inversion. Keeping the
action and charge bases, \eqref{iii:eq:dimensions} yields
\[
 \widehat\mu'=d\widehat\mu,\qquad
 \widehat\varepsilon'=d\widehat\varepsilon,\qquad
 \widehat\mu'\widehat\varepsilon'=(\widehat c')^{-2}.
\]
In the adapted chart, $\widehat\mu'=r_-/r_+$ and
$\widehat\varepsilon'=r_+/r_-$, by substitution of
$\widehat\mu=r_-^2$ and $\widehat\varepsilon=r_+^2$.

\subsubsection{Preservation of the readout under amplification}

For the refinement $\mathcal V_m=\mathcal V\otimes I_{9^m}$, the
normalized trace $\tau_m=\operatorname{Tr}/(2\cdot9^m)$ satisfies
\begin{equation}
 \exp[-2\tau_m(\log\mathcal V_m)]
 =\exp[-\log r_+-\log r_-]=\widehat c.
 \label{iii:eq:velocidad-traza-normalizada}
\end{equation}
Indeed, each eigenvalue appears $9^m$ times; that multiplicity
cancels against the denominator of the normalized trace. The ordinary
amplified determinant is $(r_+r_-)^{9^m}$. The two-sheet normalization,
equivalently expressed by \eqref{iii:eq:velocidad-traza-normalizada},
preserves the constitutive readout during amplification. This
realization preserves the enriched transport from which it originates.
